% Dot Product scaling graphs - three side by side
% Usage: % Dot Product scaling graphs - three side by side
% Usage: % Dot Product scaling graphs - three side by side
% Usage: % Dot Product scaling graphs - three side by side
% Usage: \input{figures/scaling-vecdot-graphs}

\begin{figure*}[t]
  \centering
  \begin{tikzpicture}
    \begin{groupplot}[
      group style={
        group size=3 by 1,
        horizontal sep=1.2cm,
      },
      width=0.36\textwidth,
      height=5cm,
      xlabel={$N$},
      xmode=log,
      log basis x={2},
      xmin=1.5, xmax=96,
      xtick={2,4,8,16,32,64},
      xticklabels={2,4,8,16,32,64},
      legend style={
        at={(0.5,1.02)},
        anchor=south,
        legend columns=2,
        font=\footnotesize,
      },
      grid=major,
      grid style={gray!30},
      tick label style={font=\footnotesize},
      label style={font=\small},
    ]

    % Graph 1: Qubits (linear scale to show differences clearly)
    \nextgroupplot[
      ylabel={Qubits},
      ymin=0, ymax=7500,
      title={\textbf{(a) Qubit Count}},
      title style={font=\small},
    ]
    \addplot[blue, thick, mark=square*, mark size=2pt] coordinates {
      (2, 80) (4, 144) (8, 272) (16, 528) (32, 1040) (64, 2064)
    };
    \addplot[red, thick, mark=triangle*, mark size=2pt] coordinates {
      (2, 226) (4, 436) (8, 856) (16, 1696) (32, 3376) (64, 6736)
    };
    \legend{QFT, Ripple}

    % Graph 2: Depth
    \nextgroupplot[
      ylabel={Depth},
      ymode=log,
      ymin=2500, ymax=700000,
      title={\textbf{(b) Circuit Depth}},
      title style={font=\small},
    ]
    \addplot[blue, thick, mark=square*, mark size=2pt] coordinates {
      (2, 3478) (4, 3670) (8, 4054) (16, 4822) (32, 6358) (64, 9430)
    };
    \addplot[red, thick, mark=triangle*, mark size=2pt] coordinates {
      (2, 14681) (4, 29111) (8, 57971) (16, 115691) (32, 231131) (64, 462011)
    };

    % Graph 3: T-count
    \nextgroupplot[
      ylabel={T-count},
      ymode=log,
      ymin=5000, ymax=20000000,
      title={\textbf{(c) T-count}},
      title style={font=\small},
    ]
    \addplot[blue, thick, mark=square*, mark size=2pt] coordinates {
      (2, 288754) (4, 577508) (8, 1155016) (16, 2310032) (32, 4620064) (64, 9240128)
    };
    \addplot[red, thick, mark=triangle*, mark size=2pt] coordinates {
      (2, 8910) (4, 17810) (8, 35610) (16, 71210) (32, 142410) (64, 284810)
    };

    \end{groupplot}
  \end{tikzpicture}
  \caption{Scaling behavior of dot product ($s = \sum_i a[i] \cdot b[i]$) with 8-bit integers. The ripple-carry backend requires more qubits due to ancilla allocation but achieves 32$\times$ lower T-count. QFT has lower depth due to parallel phase rotations, while ripple-carry depth grows linearly with sequential carry propagation.}
  \label{fig:scaling-vecdot}
\end{figure*}


\begin{figure*}[t]
  \centering
  \begin{tikzpicture}
    \begin{groupplot}[
      group style={
        group size=3 by 1,
        horizontal sep=1.2cm,
      },
      width=0.36\textwidth,
      height=5cm,
      xlabel={$N$},
      xmode=log,
      log basis x={2},
      xmin=1.5, xmax=96,
      xtick={2,4,8,16,32,64},
      xticklabels={2,4,8,16,32,64},
      legend style={
        at={(0.5,1.02)},
        anchor=south,
        legend columns=2,
        font=\footnotesize,
      },
      grid=major,
      grid style={gray!30},
      tick label style={font=\footnotesize},
      label style={font=\small},
    ]

    % Graph 1: Qubits (linear scale to show differences clearly)
    \nextgroupplot[
      ylabel={Qubits},
      ymin=0, ymax=7500,
      title={\textbf{(a) Qubit Count}},
      title style={font=\small},
    ]
    \addplot[blue, thick, mark=square*, mark size=2pt] coordinates {
      (2, 80) (4, 144) (8, 272) (16, 528) (32, 1040) (64, 2064)
    };
    \addplot[red, thick, mark=triangle*, mark size=2pt] coordinates {
      (2, 226) (4, 436) (8, 856) (16, 1696) (32, 3376) (64, 6736)
    };
    \legend{QFT, Ripple}

    % Graph 2: Depth
    \nextgroupplot[
      ylabel={Depth},
      ymode=log,
      ymin=2500, ymax=700000,
      title={\textbf{(b) Circuit Depth}},
      title style={font=\small},
    ]
    \addplot[blue, thick, mark=square*, mark size=2pt] coordinates {
      (2, 3478) (4, 3670) (8, 4054) (16, 4822) (32, 6358) (64, 9430)
    };
    \addplot[red, thick, mark=triangle*, mark size=2pt] coordinates {
      (2, 14681) (4, 29111) (8, 57971) (16, 115691) (32, 231131) (64, 462011)
    };

    % Graph 3: T-count
    \nextgroupplot[
      ylabel={T-count},
      ymode=log,
      ymin=5000, ymax=20000000,
      title={\textbf{(c) T-count}},
      title style={font=\small},
    ]
    \addplot[blue, thick, mark=square*, mark size=2pt] coordinates {
      (2, 288754) (4, 577508) (8, 1155016) (16, 2310032) (32, 4620064) (64, 9240128)
    };
    \addplot[red, thick, mark=triangle*, mark size=2pt] coordinates {
      (2, 8910) (4, 17810) (8, 35610) (16, 71210) (32, 142410) (64, 284810)
    };

    \end{groupplot}
  \end{tikzpicture}
  \caption{Scaling behavior of dot product ($s = \sum_i a[i] \cdot b[i]$) with 8-bit integers. The ripple-carry backend requires more qubits due to ancilla allocation but achieves 32$\times$ lower T-count. QFT has lower depth due to parallel phase rotations, while ripple-carry depth grows linearly with sequential carry propagation.}
  \label{fig:scaling-vecdot}
\end{figure*}


\begin{figure*}[t]
  \centering
  \begin{tikzpicture}
    \begin{groupplot}[
      group style={
        group size=3 by 1,
        horizontal sep=1.2cm,
      },
      width=0.36\textwidth,
      height=5cm,
      xlabel={$N$},
      xmode=log,
      log basis x={2},
      xmin=1.5, xmax=96,
      xtick={2,4,8,16,32,64},
      xticklabels={2,4,8,16,32,64},
      legend style={
        at={(0.5,1.02)},
        anchor=south,
        legend columns=2,
        font=\footnotesize,
      },
      grid=major,
      grid style={gray!30},
      tick label style={font=\footnotesize},
      label style={font=\small},
    ]

    % Graph 1: Qubits (linear scale to show differences clearly)
    \nextgroupplot[
      ylabel={Qubits},
      ymin=0, ymax=7500,
      title={\textbf{(a) Qubit Count}},
      title style={font=\small},
    ]
    \addplot[blue, thick, mark=square*, mark size=2pt] coordinates {
      (2, 80) (4, 144) (8, 272) (16, 528) (32, 1040) (64, 2064)
    };
    \addplot[red, thick, mark=triangle*, mark size=2pt] coordinates {
      (2, 226) (4, 436) (8, 856) (16, 1696) (32, 3376) (64, 6736)
    };
    \legend{QFT, Ripple}

    % Graph 2: Depth
    \nextgroupplot[
      ylabel={Depth},
      ymode=log,
      ymin=2500, ymax=700000,
      title={\textbf{(b) Circuit Depth}},
      title style={font=\small},
    ]
    \addplot[blue, thick, mark=square*, mark size=2pt] coordinates {
      (2, 3478) (4, 3670) (8, 4054) (16, 4822) (32, 6358) (64, 9430)
    };
    \addplot[red, thick, mark=triangle*, mark size=2pt] coordinates {
      (2, 14681) (4, 29111) (8, 57971) (16, 115691) (32, 231131) (64, 462011)
    };

    % Graph 3: T-count
    \nextgroupplot[
      ylabel={T-count},
      ymode=log,
      ymin=5000, ymax=20000000,
      title={\textbf{(c) T-count}},
      title style={font=\small},
    ]
    \addplot[blue, thick, mark=square*, mark size=2pt] coordinates {
      (2, 288754) (4, 577508) (8, 1155016) (16, 2310032) (32, 4620064) (64, 9240128)
    };
    \addplot[red, thick, mark=triangle*, mark size=2pt] coordinates {
      (2, 8910) (4, 17810) (8, 35610) (16, 71210) (32, 142410) (64, 284810)
    };

    \end{groupplot}
  \end{tikzpicture}
  \caption{Scaling behavior of dot product ($s = \sum_i a[i] \cdot b[i]$) with 8-bit integers. The ripple-carry backend requires more qubits due to ancilla allocation but achieves 32$\times$ lower T-count. QFT has lower depth due to parallel phase rotations, while ripple-carry depth grows linearly with sequential carry propagation.}
  \label{fig:scaling-vecdot}
\end{figure*}


\begin{figure*}[t]
  \centering
  \begin{tikzpicture}
    \begin{groupplot}[
      group style={
        group size=3 by 1,
        horizontal sep=1.2cm,
      },
      width=0.36\textwidth,
      height=5cm,
      xlabel={$N$},
      xmode=log,
      log basis x={2},
      xmin=1.5, xmax=96,
      xtick={2,4,8,16,32,64},
      xticklabels={2,4,8,16,32,64},
      legend style={
        at={(0.5,1.02)},
        anchor=south,
        legend columns=2,
        font=\footnotesize,
      },
      grid=major,
      grid style={gray!30},
      tick label style={font=\footnotesize},
      label style={font=\small},
    ]

    % Graph 1: Qubits (linear scale to show differences clearly)
    \nextgroupplot[
      ylabel={Qubits},
      ymin=0, ymax=7500,
      title={\textbf{(a) Qubit Count}},
      title style={font=\small},
    ]
    \addplot[blue, thick, mark=square*, mark size=2pt] coordinates {
      (2, 80) (4, 144) (8, 272) (16, 528) (32, 1040) (64, 2064)
    };
    \addplot[red, thick, mark=triangle*, mark size=2pt] coordinates {
      (2, 226) (4, 436) (8, 856) (16, 1696) (32, 3376) (64, 6736)
    };
    \legend{QFT, Ripple}

    % Graph 2: Depth
    \nextgroupplot[
      ylabel={Depth},
      ymode=log,
      ymin=2500, ymax=700000,
      title={\textbf{(b) Circuit Depth}},
      title style={font=\small},
    ]
    \addplot[blue, thick, mark=square*, mark size=2pt] coordinates {
      (2, 3478) (4, 3670) (8, 4054) (16, 4822) (32, 6358) (64, 9430)
    };
    \addplot[red, thick, mark=triangle*, mark size=2pt] coordinates {
      (2, 14681) (4, 29111) (8, 57971) (16, 115691) (32, 231131) (64, 462011)
    };

    % Graph 3: T-count
    \nextgroupplot[
      ylabel={T-count},
      ymode=log,
      ymin=5000, ymax=20000000,
      title={\textbf{(c) T-count}},
      title style={font=\small},
    ]
    \addplot[blue, thick, mark=square*, mark size=2pt] coordinates {
      (2, 288754) (4, 577508) (8, 1155016) (16, 2310032) (32, 4620064) (64, 9240128)
    };
    \addplot[red, thick, mark=triangle*, mark size=2pt] coordinates {
      (2, 8910) (4, 17810) (8, 35610) (16, 71210) (32, 142410) (64, 284810)
    };

    \end{groupplot}
  \end{tikzpicture}
  \caption{Scaling behavior of dot product ($s = \sum_i a[i] \cdot b[i]$) with 8-bit integers. The ripple-carry backend requires more qubits due to ancilla allocation but achieves 32$\times$ lower T-count. QFT has lower depth due to parallel phase rotations, while ripple-carry depth grows linearly with sequential carry propagation.}
  \label{fig:scaling-vecdot}
\end{figure*}
